\section{Rank-one harmonic dependence and depth-two closure}
\label{dhr:sec:depth}

\subsection{A theorem beyond consecutive factorial weights}
The next result is algebraic and does not require the Gamma product.
Let $w(x)\in\Q(x)$ have only simple poles, all at negative integers,
and suppose
\begin{equation}\label{dhr:eq:rational-weight-assumption}
 w(n)=\lambda n^{-(2r-1)}+O(n^{-2r}),\qquad
 \lambda\in\Q\setminus\{0\}.
\end{equation}
There is no positivity assumption. The pole locations need not be
consecutive. Define
\begin{equation}\label{dhr:eq:general-cutoff}
 S_{w,N}=\sum_{0\leq n_1<\cdots<n_r\leq N}
              \Delta(\bm n)^2\prod_iw(n_i),\qquad S_{w,-1}=0.
\end{equation}

\begin{theorem}[Rank-one harmonic dependence]\label{dhr:thm:rank-one}
There exist rational functions $P_w(N),Q_w(N)\in\Q(N)$ such that
for every integer $N\geq0$,
\begin{equation}\label{dhr:eq:rank-one-cutoff}
 \boxed{S_{w,N}=P_w(N)H_{N+1}+Q_w(N).}
\end{equation}
Both rational functions are bounded at infinity. At each finite pole,
$P_w$ has order at most $r-1$ and $Q_w$ has order at most $r$.
The possible poles are negative integers at most $-2$.
For $r=1$, $P_w=\lambda$.
\end{theorem}

\begin{proof}
For $0\leq k\leq2r-2$, partial fractions give
\begin{equation}\label{dhr:eq:moment-residues}
 n^kw(n)=\sum_{d\geq1}\frac{c_{k,d}}{n+d},\qquad
 \sum_dc_{k,d}=\lambda\ind_{k=2r-2}.
\end{equation}
Only finitely many coefficients are nonzero. There is no polynomial
part because $n^kw(n)$ is proper. The sum of residues is the coefficient
of $n^{-1}$ at infinity, proving the second equality.

Thus the finite moment $M_k(N)=\sum_{n=0}^Nn^kw(n)$ equals
\begin{align}
 M_k(N)&=R_k(N)+\lambda\ind_{k=2r-2}H_{N+1},\notag\\
 R_k(N)&=\sum_dc_{k,d}\left\{
       \sum_{\ell=2}^{d}\frac1{N+\ell}-H_{d-1}\right\}.
 \label{dhr:eq:finite-moments}
\end{align}
Every $R_k$ is a bounded rational function with simple poles only
at the indicated negative integers. Finite Cauchy--Binet gives
\begin{equation}\label{dhr:eq:hankel-cutoff}
 S_{w,N}=\det_{0\leq i,j<r}M_{i+j}(N).
\end{equation}
Among these entries, $i+j=2r-2$ occurs only at $i=j=r-1$.
Therefore
\begin{equation}\label{dhr:eq:rank-one-matrix}
 (M_{i+j})_{i,j=0}^{r-1}
 =(R_{i+j})_{i,j=0}^{r-1}
       +\lambda H_{N+1}\bm e_{r-1}\bm e_{r-1}^{\mathsf T}.
\end{equation}
The determinant is affine in the harmonic number, with
\begin{equation}\label{dhr:eq:PQ-det}
 P_w=\lambda\det_{0\leq i,j<r-1}R_{i+j},\qquad
 Q_w=\det_{0\leq i,j<r}R_{i+j}.
\end{equation}
The empty-determinant convention covers $r=1$. Products of $r-1$
or $r$ simple-pole entries prove the pole-order bounds, and boundedness
of every entry proves boundedness at infinity. This argument proves
\eqref{dhr:eq:rank-one-cutoff} at every cutoff, including $N<r-1$
where the determinant and configuration sum both vanish.
\end{proof}

\begin{remark}[The source of the depth collapse]
A generic determinant whose entries are affine in $H_{N+1}$ has
harmonic degree as large as $r$. That estimate is much too weak here.
The decay condition \eqref{dhr:eq:rational-weight-assumption} forces
the residue sum to vanish in every moment except the highest one.
Consequently the harmonic matrix has rank one \emph{and a single
nonzero entry}. This is the mechanism, not cancellation observed in
a few low-rank examples.
\end{remark}

\subsection{An explicit partial-fraction transform}
For later use put
\[
 T_{p,j}(z)=\sum_{m\geq1}\frac{z^m}{(m+j)^p},\qquad
 U_{p,j}(z)=\sum_{m\geq1}\frac{z^mH_m}{(m+j)^p},
 \qquad p\geq1,\ j\geq0.
\]
Shifting the summation index gives
\begin{equation}\label{dhr:eq:T-transform}
 T_{p,j}(z)=z^{-j}\left\{\Li_p(z)-
                    \sum_{k=1}^j\frac{z^k}{k^p}\right\}.
\end{equation}
For the harmonic sums the precise shifted formula is
\begin{align}
 U_{p,0}(z)&=\Li_{p,1}(z,1)+\Li_{p+1}(z),
 \label{dhr:eq:U-zero}\\
 U_{p,j}(z)&=z^{-j}\bigg\{\Li_{p,1}(z,1)
 -\sum_{k=1}^j\frac{z^kH_{k-1}}{k^p}
 -\sum_{\ell=1}^{j-1}\sum_{k=j+1}^{\infty}
             \frac{z^k}{k^p(k-\ell)}\bigg\},\quad j\geq1.
 \label{dhr:eq:U-shift}
\end{align}
For \eqref{dhr:eq:U-shift}, write $k=m+j$ and use
$H_m=H_{k-1}-\sum_{\ell=1}^{j-1}(k-\ell)^{-1}$.
The finite correction is $H_{k-1}$, not $H_k$.
Each last rational tail has a partial-fraction decomposition in
$k$ and $k-\ell$, hence involves only ordinary polylogarithms and
finite sums. Its starting index $j+1>\ell$ avoids the apparent pole.
Equations \eqref{dhr:eq:T-transform}--\eqref{dhr:eq:U-shift} are
explicit terminating reduction rules, not formal existence statements.

\begin{theorem}[All-particle depth-two closure]\label{dhr:thm:closure}
Under the assumptions of Theorem~\ref{dhr:thm:rank-one}, let
$K_{w,N}=S_{w,N}-S_{w,N-1}$ and
\begin{equation}\label{dhr:eq:general-generating}
 \mathcal G_w(z)=\sum_{N\geq0}K_{w,N}z^{N+1},\qquad |z|<1.
\end{equation}
Then $\mathcal G_w$ belongs to the $\Q(z)$-linear span of
\begin{equation}\label{dhr:eq:closure-space}
 1,\quad \Li_1(z),\ldots,\Li_r(z),\quad
 \Li_{1,1}(z,1),\ldots,\Li_{r-1,1}(z,1).
\end{equation}
For $r=1$ the depth-two part of this list is empty. Thus the depth
upper bound is two, independently of $r$, and the polylogarithmic
weight upper bound is $r$.
\end{theorem}

\begin{proof}
Set $m=N+1$. By \eqref{dhr:eq:rank-one-cutoff},
\begin{align}
 K_{w,N}={}&\{P_w(N)-P_w(N-1)\}H_{N+1}\notag\\
          &+Q_w(N)-Q_w(N-1)+\frac{P_w(N-1)}{N+1}.
 \label{dhr:eq:mass-difference}
\end{align}
The harmonic coefficient is a proper rational function of $m$,
with poles only at nonpositive integers and pole orders at most $r-1$.
It is zero for $r=1$. The nonharmonic coefficient is also proper,
with pole orders at most $r$: division by $m$ can increase the
order of $P_w(N-1)$ at zero only from at most $r-1$ to at most $r$.
All other poles already have the stated bound. The expressions at
$N=-1$ used in \eqref{dhr:eq:mass-difference} are well defined;
\eqref{dhr:eq:finite-moments} gives the empty moments there.

Decompose both coefficients into terms $(m+j)^{-p}$ with $j\geq0$.
For the nonharmonic part apply \eqref{dhr:eq:T-transform}, which
uses indices at most $r$. For the harmonic part apply
\eqref{dhr:eq:U-zero} and \eqref{dhr:eq:U-shift}. The resulting depth-two
index is at most $r-1$; any additional ordinary polylogarithm has
index at most $r$. This proves \eqref{dhr:eq:closure-space}.
All the series converge absolutely in the open disk, so the shifts
and finite partial fractions are legitimate.
\end{proof}

The equivalent identity
\begin{equation}\label{dhr:eq:Abel-cutoff}
 \mathcal G_w(z)=(1-z)\sum_{N\geq0}S_{w,N}z^{N+1}
\end{equation}
connects the cutoff and Abel descriptions. It is useful for checking
constants but is not needed to guess the depth bound.

\subsection{Derivatives, Euler primitives, and endpoint zeta reduction}
Let $\Theta=z\,d/dz$ and normalize $\Theta^{-1}$ by integration from
zero with the kernel $dt/t$. Then
\begin{equation}\label{dhr:eq:Euler-primitive}
 \Theta^{-q}\mathcal G_w(z)=
 \sum_{N\geq0}\frac{K_{w,N}}{(N+1)^q}z^{N+1},\qquad q\geq0.
\end{equation}
This also equals, for $q\geq1$,
\begin{equation}\label{dhr:eq:Euler-integral}
 \frac1{(q-1)!}\int_0^z\log^{q-1}(z/t)
                       \mathcal G_w(t)\frac{\mathrm dt}{t},
\end{equation}
initially along $0<t<z<1$, then as the analytic primitive of the
power series.

\begin{corollary}[Primitive closure and ordinary endpoint evaluations]
\label{dhr:cor:endpoint}
For every $q\geq0$, the function \eqref{dhr:eq:Euler-primitive} has
polylogarithmic depth at most two and weight at most $r+q$.
For $q\geq1$ its endpoint
\begin{equation}\label{dhr:eq:I-general}
 I_{w,q}=\sum_{N\geq0}\frac{K_{w,N}}{(N+1)^q}
\end{equation}
is an ordinary absolutely convergent sum and lies in the rational
linear span of
\[
 1,\qquad \zeta(k),\qquad \zeta(a)\zeta(b),
 \quad 2\leq k\leq r+q,\quad a,b\geq2,\quad a+b\leq r+q.
\]
No product of more than two zeta factors is needed for this endpoint
family. Ordinary derivatives and anchored ordinary antiderivatives
also retain depth at most two, although their rational prefactors
and exponent shifts change.
\end{corollary}

\begin{proof}
Divide the two proper rational coefficients in
\eqref{dhr:eq:mass-difference} by $m^q$. Their maximum pole orders
become at most $r-1+q$ in the harmonic part and $r+q$ in the other
part. The preceding proof therefore gives the asserted weight and
depth bounds.

Bounded rational functions have expansions at infinity, so
\eqref{dhr:eq:mass-difference} also gives
$K_{w,N}=P_w(\infty)/(N+1)+O(\log(N+2)/(N+1)^2)$.
This proves absolute convergence of \eqref{dhr:eq:I-general}.
The configuration sum itself is absolutely convergent after inserting
the maximum weight $(n_r+1)^{-q}$, by the same absolute
largest-index estimate used in Section~\ref{dhr:sec:resonance}.

The reduction rules give Laurent-polynomial, hence locally analytic
near $z=1$, coefficients of the polylogarithms. Only $\Li_1(z)$
and $\Li_{1,1}(z,1)=\tfrac12\log^2(1-z)$ diverge there. Convergence
forces their coefficients to vanish at one: a nonzero quadratic
logarithm cannot cancel a linear logarithm or a bounded term, and
a nonzero linear logarithm cannot cancel bounded terms. Those
vanishing-coefficient terms tend to zero. The other endpoint values
are $\zeta(p)$ and $\zeta(p,1)$. Lemma~\ref{dhr:lem:Euler} below
reduces the latter to the stated span.

For the last assertion, power-series differentiation multiplies
coefficients by affine factors in the index and shifts the exponent;
ordinary integration divides by such factors and shifts it oppositely.
A finite shift of $H_m$ changes it by a rational function. Polynomial
parts in the index, when they arise, sum to rational functions or
rational multiples of $\Li_1(z)$. Partial fractions of the remaining
proper parts again use \eqref{dhr:eq:T-transform}--\eqref{dhr:eq:U-shift},
with finite initial terms removed where needed. None of these operations
requires a third nested index.
\end{proof}

\begin{lemma}[Euler's height-one reduction]\label{dhr:lem:Euler}
For every integer $p\geq2$,
\begin{equation}\label{dhr:eq:Euler-zeta}
 \zeta(p,1)=\frac p2\zeta(p+1)
 -\frac12\sum_{j=1}^{p-2}\zeta(p-j)\zeta(j+1).
\end{equation}
\end{lemma}
\begin{proof}
We include a beta-integral proof of this classical identity so the
endpoint theorem has no unproved multiple-zeta reduction step.
For $|x|<1$, expand Euler's beta integral \cite{dlmf:beta}:
\begin{align*}
 B(1-x,y)-B(1,y)
 &=\int_0^1(t^{-x}-1)(1-t)^{y-1}\mathrm dt\\
 &=-\gamma-\psi(1-x)\\
 &\quad+\frac y2\big\{[\gamma+\psi(1-x)]^2
                          +\zeta(2)-\psi'(1-x)\big\}+O(y^2).
\end{align*}
The $y$-derivative at zero can be taken under the integral: the
factor $t^{-x}-1$ cancels its endpoint singularity at $t=1$.
Using $\log(1-t)/(1-t)=-\sum_{n\geq1}H_nt^n$ gives this derivative
as
\[
 -\sum_{n\geq1}H_n\left\{\frac1{n+1-x}-\frac1{n+1}\right\}
 =-\sum_{p\geq2}\zeta(p,1)x^{p-1}.
\]
The interchanges hold first for real $0<x<1$ by absolute convergence,
and then throughout the disk by holomorphy. Put
$f(x)=-\psi(1-x)-\gamma=\sum_{k\geq2}\zeta(k)x^{k-1}$.
Comparison with the beta expansion yields
\[
 \sum_{p\geq2}\zeta(p,1)x^{p-1}
 =\tfrac12\{f'(x)-\zeta(2)-f(x)^2\}.
\]
Equating coefficients of $x^{p-1}$ proves
\eqref{dhr:eq:Euler-zeta}.
\end{proof}

\subsection{A common translate and cyclotomic coordinates}
The rank-one mechanism survives translation of the entire simple-pole
grid. Define
\[
 H_n(a)=\sum_{j=0}^{n-1}\frac1{j+a},\qquad H_0(a)=0,\qquad a>0.
\]
Suppose all poles of $w(n)$ are simple and lie at $-a-d$ with
integers $d\geq0$, and the critical decay
\eqref{dhr:eq:rational-weight-assumption} holds. Use the monic moment
basis $(n+a)^k$. For each pole,
\[
 \sum_{n=0}^N\frac1{n+a+d}
 =H_{N+1}(a)+\sum_{\ell=1}^d\frac1{N+a+\ell}
                 -\sum_{\ell=0}^{d-1}\frac1{a+\ell}.
\]
The sum of residues still vanishes except in the top moment.
The proof of Theorem~\ref{dhr:thm:rank-one} therefore gives
\begin{equation}\label{dhr:eq:translated-cutoff}
 S_{w,N}=P_w(N,a)H_{N+1}(a)+Q_w(N,a),
\end{equation}
with the same pole-order bounds. For a rational weight depending
rationally on $a$, the coefficients are rational in $N,a$.

For rational $a=p/q\in(0,1]$ in lowest terms, put
$\omega=e^{2\pi i/q}$. The exact root-of-unity filters are
\begin{align}
 \sum_{n\geq0}\frac{z^{n+a}}{(n+a)^s}
 &=q^{s-1}\sum_{j=0}^{q-1}\omega^{-pj}
                       \Li_s(\omega^jz^{1/q}),
 \label{dhr:eq:single-filter}\\
 \sum_{n\geq0}\frac{z^{n+a}H_n(a)}{(n+a)^s}
 &=q^{s-1}\sum_{j,k=0}^{q-1}\omega^{-p(j+k)}
             \Li_{s,1}(\omega^jz^{1/q},\omega^k).
 \label{dhr:eq:double-filter}
\end{align}
These formulas hold for integers $s\geq1$ and $0<z<1$, and on any
common continuation with a fixed $q$th-root branch. To prove them,
write the top index as $qn+p$ and the inner index as $qm+p$.
The powers of $q$ from the denominators are $q^s$ and $q^{s+1}$;
the one and two residue filters divide by $q$ and $q^2$,
respectively. The strict inner inequality is preserved.

Finite index shifts and partial fractions, exactly as in
\eqref{dhr:eq:U-shift}, now prove depth at most two in level-$q$
colored polylogarithms for the generating function with exponent
$z^{N+a}$. The same conclusion holds for primitives that divide
its coefficients by $(N+a)^h$. No ordinary-zeta-only endpoint
reduction is asserted for general level $q$.

The shifted canonical weight $w(n)=1/(n+a)_{2r-1}$ is also covered
by the Gamma product: take the parameters there to be $(a,c)=(a,1)$
and $b_0=a+2r-1$. For $r=2$ its residue and cutoff constant are
\begin{equation}\label{dhr:eq:translated-cutoff-constant}
 A_0=\frac1{2a(a+1)},\qquad
 C_a=\frac{-2-\psi(a+1)}{2a(a+1)}.
\end{equation}
For example, at $a=1/2$,
\begin{equation}\label{dhr:eq:half-cutoff}
 \lim_{N\to\infty}\left[
 \sum_{0\leq m<n\leq N}
 \frac{(n-m)^2}{(m+\tfrac12)_3(n+\tfrac12)_3}
       -\frac23\log N\right]
 =\frac23\{\gamma+2\log2-4\}.
\end{equation}
This gives a direct digamma/cyclotomic bridge without changing the
finite-cutoff depth mechanism.
